\begin{corollary}[A tradeoff up to exponent constant three]
\label{tou:degree-tradeoff}
In either two-orbit construction, one may instead take
$K=dD-s$ for any $1\le s<d$. The nearby count and uniqueness
conclusions persist, with
\[
 \eta=\frac{2d+s}{n},\qquad
 \dist(f,C)=1-\rho-s/n
          =\theta+\frac{2d}{2d+s}\eta.
\]
The numerical prescription is violated exponentially for every
fixed $c_2<2+s/d$. Thus any fixed $0<c_2<3$ is compatible with
unique nearby witnesses and separation greater than two thirds of
the gap. Efficient witness recovery remains available in the
integer construction; the function-field construction retains its
logarithmic length and existence guarantees.
\end{corollary}
\begin{proof}
All constructed differences $w-H_I$ have degree at most $dD-d<K$.
For any codeword, the monic difference from $w$ still has degree
$dD=K+s$ and zero penultimate coefficient. The same root-sum
argument classifies every codeword with $dD+2d$ agreements.
The far word has exactly $dD$ agreements, so both displayed distances
follow. For an exact rational rate $\rho=a/b$, choose $n=bu$ with
$au\equiv-s\pmod d$; this is possible when $d\nmid a$.
The support count still has logarithm $nH_2(\rho)/d+O(\log n)$,
which gives the threshold $c_2<(2d+s)/d$.
Take $s=d-1$ and then a sufficiently large prime $d$ for the final
claim. In the logarithmic-length version the strict Elias condition
still holds; in the integer version choose the splitting prime
sufficiently large.
\end{proof}

\begin{remark}[The dimension cutoff in the integer construction]
For $D\ge2$, lowering $K$ further to $dD-d$ leaves at most one nearby
parameter at the same agreement threshold $dD+2d$.
Indeed the sums $\sum_{i\in I}R_i$ of distinct $D$-supports are
distinct. To see this over the rationals, write
$R_i=q+(q-1)/(q^i-1)$ and take the least index in the symmetric
difference. Its nonconstant term exceeds the sum of all later terms,
since $(q-1)^2>q$. Multiplication by $W$ clears denominators, and
the nonzero difference has absolute value at most $Wn(q+1)<p$;
thus distinctness survives reduction.
The coefficient of $X^{dD-d}$ in $H_I$ is $-\sum_{i\in I}R_i$.
For $w-H_I$ to have degree below $dD-d$, this sum must equal the
reference support's sum, forcing that support alone.
The all-codeword classification excludes any other nearby witnesses.
This is a limitation of further dimension reduction in this construction,
not an upper bound for other constructions.
\end{remark}
